\section*{Hoofdstuk \ref{ch:b3:complete} --- Volledige ruimten:
Baire, Ascoli, Stone--Weierstrass}

\begin{solution}{exo:b3:complete:1}
(a) Voor $m \geq n$ verdwijnt $f_n - f_m$ buiten een interval van
lengte $\frac1n$ en is ze begrensd door $1$: $\norm{f_n - f_m}_1
\leq \frac1n \to 0$, dus Cauchy. Was $f_n \to f$ in
$\norm\cdot_1$ met $f$ \omterm{def:b3:topology:continuity}{continu}, dan is voor een vaste $\alpha <
\frac12$: $\int_0^{\alpha}\abs f = \lim \int_0^\alpha\abs{f - f_n
+ f_n} \leq \lim\bigl(\norm{f - f_n}_1 + 0\bigr) = 0$ (want $f_n
\equiv 0$ daar voor grote $n$), dus $f \equiv 0$ op $[0,
\frac12)$ (\omterm{def:b3:topology:continuity}{continuïteit}); en net zo $f \equiv 1$ op $(\frac12,
1]$: geen \omterm{def:b3:topology:continuity}{continue} functie doet beide. De ruimte is dus onvolledig
--- de \omterm{thm:b3:complete:completion}{vervollediging} is $L^1$, gebouwd in
\cref{ch:b3:lp}.

(b) Zij $(x_n)$ Cauchy; kies $n_1 < n_2 < \cdots$ met
$\norm{x_{n_{k+1}} - x_{n_k}} \leq 2^{-k}$. De reeks $\sum_k
(x_{n_{k+1}} - x_{n_k})$ convergeert absoluut en dus; haar
partiële sommen zijn $x_{n_{k+1}} - x_{n_1}$, dus convergeert
$(x_{n_k})$, en een Cauchyrij met een convergente deelrij
convergeert.
\end{solution}

\begin{solution}{exo:b3:complete:2}
(a) Stel dat $E$ \omterm{def:b3:complete:complete}{volledig} is met \omterm{def:b3:galois:algebraic}{algebraïsche} basis
$(e_n)_{n\in\N}$, en stel $F_n = \operatorname{Vect}(e_1, \dots,
e_n)$: gesloten (eindigdimensionale deelruimten zijn \omterm{def:b3:complete:complete}{volledig} en
dus gesloten --- bachelorjaar 2), met leeg inwendige: was $B(x, r)
\subseteq F_n$, neem dan $v \notin F_n$; dan is $x + \frac{r}{2\norm
v}v \in B(x,r) \setminus F_n$, absurd. Maar $E = \bigcup F_n$
(elke vector is een eindige combinatie): in strijd met Baire
(\cref{thm:b3:complete:baire}). De ruimte $\R[X]$ heeft de
aftelbare basis $(X^n)$, dus maakt geen enkele norm haar
\omterm{def:b3:complete:complete}{volledig}.

(b) Leg $\delta > 0$ en een niet-lege open bal $B$ vast; we zoeken
in $B$ een punt van $\Omega_\delta = \{x : \operatorname{osc}_xf <
\delta\}$, waarbij $\operatorname{osc}_xf = \inf_{V \ni
x}\operatorname{diam} f(V)$. ($\Omega_\delta$ is open: is
$\operatorname{diam}f(V) < \delta$ voor een open $V \ni x$, dan
heeft elke $y \in V$ oscillatie $< \delta$.) De verzamelingen
\[
F_N = \bigl\{x : \abs{f_n(x) - f_m(x)} \leq \tfrac\delta3\
\ \forall\, n, m \geq N \bigr\}
\]
zijn gesloten (doorsneden van originelen van gesloten
verzamelingen) en overdekken $\R$ (de puntsgewijze convergentie
maakt $(f_n(x))$ Cauchy). Baire toegepast binnen de volledige
$\bar B$: een zekere $F_N \cap \bar B$ bevat een bal $B' = B(x_0,
\rho)$. Laat $m \to \infty$ gaan: $\abs{f_N - f} \leq
\frac\delta3$ op $B'$. Krimp wegens de \omterm{def:b3:topology:continuity}{continuïteit} van $f_N$ in
$x_0$ tot een $B'' \ni x_0$ waar $\abs{f_N - f_N(x_0)} \leq
\frac\delta3$; dan is voor $x \in B''$
\[
\abs{f(x) - f(x_0)} \leq \abs{f - f_N}(x) + \abs{f_N(x) -
f_N(x_0)} + \abs{f_N - f}(x_0) \leq \delta:
\]
dus $\operatorname{diam} f(B'') \leq 2\delta$ en $x_0 \in
\Omega_{3\delta} \cap B$. Elke $\Omega_\delta$ is dus open en
\omterm{def:b3:topology:interior}{dicht}, en $\bigcap_k\Omega_{1/k}$ --- de verzameling
continuïteitspunten --- is \omterm{def:b3:topology:interior}{dicht} volgens Baire.
\end{solution}

\begin{solution}{exo:b3:complete:3}
(a) $f(x) = \frac12(x + \frac ax)$ beeldt $[\sqrt a, \infty)$ op
zichzelf af (ongelijkheid van rekenkundig en meetkundig
gemiddelde: $f(x) \geq \sqrt{x \cdot \frac ax} = \sqrt a$), en
daar is $f'(x) = \frac12(1 - \frac a{x^2}) \in [0, \frac12)$: een
$\frac12$-lipschitz contractie van een gesloten (dus volledige)
verzameling. Vast punt: $x = f(x) \iff x^2 = a$, dus $x^* = \sqrt
a$. Snelheid (\cref{thm:b3:complete:banach}): $d(x_n, \sqrt a)
\leq 2^{-n+1}\,d(x_1, x_0)$. Voor $a = 2$ en $x_0 = 2$ is $d(x_1,
x_0) = \frac12$, dus garandeert $n = 40$ dat $2^{-40} <
10^{-12}$. (In werkelijkheid convergeert de methode van Newton
kwadratisch: een handvol iteraties volstaat; de contractieschatting
is pessimistisch maar gratis.)

(b) $f_m(x) = m + e\sin x$ is $e$-lipschitz met $e < 1$ op het
volledige $\R$: dus precies één vast punt $x(m)$. Voor twee
parameters:
\[
\abs{x(m) - x(m')} = \abs{f_m(x(m)) - f_{m'}(x(m'))}
\leq \abs{m - m'} + e\abs{x(m) - x(m')},
\]
dus $\abs{x(m) - x(m')} \leq \frac{\abs{m - m'}}{1 - e}$: zelfs
lipschitz in $m$.
\end{solution}

\begin{solution}{exo:b3:complete:4}
(a) Zij $x^*$ het unieke vaste punt van $f^p$. Dan is
$f^p(f(x^*)) = f(f^p(x^*)) = f(x^*)$: $f(x^*)$ is een vast punt
van $f^p$, dus $f(x^*) = x^*$. Een vast punt van $f$ is er een van
$f^p$: de eenduidigheid draagt over. Voor $T$: per inductie is
$\abs{T^pf(x)} \leq \norm f_\infty x^p/p!$ (elke integratie voegt
een factor $\frac xk$ toe), dus $\norm{T^p} \leq a^p/p!$. De
afbeelding $S(u) = g + \lambda Tu$ voldoet aan $S^p(u) - S^p(v) =
\lambda^pT^p(u - v)$, met norm $\leq \abs\lambda^pa^p/p!\,\norm{u
- v} \to 0$: een zekere $S^p$ is een contractie, en $S$ heeft
precies één vast punt: de vergelijking van Volterra $u = g +
\lambda Tu$ is voor \emph{elke} $\lambda$ eenduidig \omterm{def:b3:groups:derived}{oplosbaar}.

(b) $\varphi(x) = d(x, f(x))$ is \omterm{def:b3:topology:continuity}{continu} op de \omterm{def:b3:topology:compact}{compacte} $K$: ze
bereikt haar minimum in een zekere $x_0$. Was $f(x_0) \neq x_0$,
dan is $\varphi(f(x_0)) = d(f(x_0), f^2(x_0)) < d(x_0, f(x_0)) =
\min\varphi$: absurd. Eenduidigheid: twee vaste punten $x \ne y$
zouden $d(x,y) = d(f(x), f(y)) < d(x,y)$ geven. Zonder
\omterm{def:b3:topology:compact}{compactheid}: $f(x) = x + \frac1x$ op $[1, \infty)$ voldoet voor $x
< y$ aan $f(y) - f(x) = (y - x)\bigl(1 - \frac1{xy}\bigr) < y -
x$, en toch is altijd $f(x) > x$: geen vast punt --- een strikte
afname van de afstanden is zwakker dan een uniforme
contractiefactor.
\end{solution}

\begin{solution}{exo:b3:complete:5}
(a) De verzameling $A = \{g = h\}$ is het origineel van de
diagonaal $\Delta_Y$ onder de \omterm{def:b3:topology:continuity}{continue} $x \mapsto (g(x), h(x))$;
en $\Delta_Y$ is gesloten omdat $Y$ \omterm{def:b3:topology:hausdorff}{Hausdorff} is (voor $y_1 \neq
y_2$ geven disjuncte open \omterm{def:b3:topology:topology}{omgevingen} een open blok rond $(y_1,
y_2)$ dat de diagonaal mijdt, zodat het complement van $\Delta_Y$
open is): dus is $A$ gesloten, bevat het een \omterm{def:b3:topology:interior}{dichte verzameling} en
is het gelijk aan $X$. Gebruikt: de eenduidigheid in
\cref{thm:b3:complete:extension}, en daarmee in de eenduidigheid
van de \omterm{thm:b3:complete:completion}{vervollediging}.

(b) $f$ is een isometrie en dus uniform \omterm{def:b3:topology:continuity}{continu}: ze zet zich voort
tot $F\colon X \to Y$ (\cref{thm:b3:complete:extension}), nog
steeds isometrisch (de betrekking $d(F(x), F(x')) = d(x, x')$
geldt op een \omterm{def:b3:topology:interior}{dichte verzameling} paren en beide leden zijn
\omterm{def:b3:topology:continuity}{continu}). Net zo zet $f^{-1}\colon f(D) \to X$ zich voort tot $G
\colon Y \to X$. De samenstelling $G \circ F$ is \omterm{def:b3:topology:continuity}{continu} en houdt
de \omterm{def:b3:topology:interior}{dichte} $D$ vast: ze is $\mathrm{id}_X$ (onderdeel (a)); en
symmetrisch is $F \circ G = \mathrm{id}_Y$. Dus is $F$ een
isometrische bijectie. Eenduidigheid van de \omterm{thm:b3:complete:completion}{vervollediging}: pas
dit toe op $D = X$, \omterm{def:b3:topology:interior}{dicht} gelegen in twee \omterm{thm:b3:complete:completion}{vervolledigingen}.
\end{solution}

\begin{solution}{exo:b3:complete:6}
$\{\sin(nx)\}$: puntsgewijs begrensd door $1$; \emph{niet}
\omterm{def:b3:complete:equicontinuous}{equicontinu}: in $x = 0$ is $\sin(n\cdot\frac{\pi}{2n}) = 1$ met
$\frac\pi{2n} \to 0$, wat elke gemeenschappelijke $\delta$ voor
$\varepsilon = \frac12$ schendt. Niet relatief \omterm{def:b3:topology:compact}{compact} (de
noodzakelijkheid in Ascoli,
\cref{thm:b3:complete:ascoli}).

$\{x^n\}$: begrensd; niet \omterm{def:b3:complete:equicontinuous}{equicontinu} in $1$: $1 - (1 - \delta)^n
\to 1$ als $n \to \infty$ voor vaste $\delta$. Niet relatief
\omterm{def:b3:topology:compact}{compact} --- en consistent daarmee is haar puntsgewijze limiet
discontinu, zodat geen enkele deelrij uniform convergeert.

$\{\norm f_\infty \leq 1, \norm{f'}_\infty \leq 1\}$: de
middelwaardeongelijkheid maakt de familie $1$-lipschitz en dus
\omterm{def:b3:complete:equicontinuous}{equicontinu}; begrensd: relatief \omterm{def:b3:topology:compact}{compact} volgens Ascoli. (Niet
\omterm{def:b3:topology:compact}{compact}: ze is niet gesloten --- uniforme limieten hoeven niet
$\mathcal C^1$ te zijn; haar \omterm{def:b3:topology:interior}{afsluiting} bestaat uit de
$1$-lipschitz functies met norm $\leq 1$.)

$\{\operatorname{Lip}(f) \leq 1, f(0) = 0\}$: \omterm{def:b3:complete:equicontinuous}{equicontinu};
puntsgewijs begrensd ($\abs{f(x)} \leq x \leq 1$); en gesloten
onder uniforme limieten (de lipschitzongelijkheid en de waarde in
$0$ gaan over op de limiet): \omterm{def:b3:topology:compact}{compact}.
\end{solution}

\begin{solution}{exo:b3:complete:7}
(a) Voor $\norm f_\infty \leq 1$ is $\abs{Tf(x)} \leq
\norm k_\infty$, en
\[
\abs{Tf(x) - Tf(x')} \leq \int_0^1\abs{k(x,y) -
k(x',y)}\,\dd y \leq \omega_k\bigl(\abs{x - x'}\bigr),
\]
waarbij $\omega_k$ een modulus van uniforme \omterm{def:b3:topology:continuity}{continuïteit} van $k$
op het \omterm{def:b3:topology:compact}{compacte} vierkant is (Heine): het beeld van de eenheidsbal
is uniform begrensd en \omterm{def:b3:complete:equicontinuous}{equicontinu}, en dus relatief \omterm{def:b3:topology:compact}{compact}
(Ascoli). Wegens de lineariteit heeft elke begrensde verzameling
een relatief \omterm{def:b3:topology:compact}{compact} beeld: $T$ is een \omterm{def:b3:topology:compact}{compacte} operator.

(b) De uitspraak over deelrijen is de definitie van relatieve
\omterm{def:b3:topology:compact}{compactheid}, toegepast op $\{Tf_n\}$. Was $T$ bijectief met
\omterm{def:b3:topology:continuity}{continue} inverse, dan is $T(B(0,1)) \supseteq B(0, \varepsilon)$
voor zekere $\varepsilon > 0$ ($T^{-1}$ is \omterm{def:b3:topology:continuity}{continu} in $0$); en dan
was de gesloten bal $\bar B(0,\varepsilon)$, een gesloten
deelverzameling van de \omterm{def:b3:topology:compact}{compacte} $\overline{T(B(0,1))}$, \omterm{def:b3:topology:compact}{compact}
--- onmogelijk in het oneindigdimensionale $\mathcal C(\intcc01)$
volgens de stelling van Riesz (bachelorjaar 2).
\end{solution}

\begin{solution}{exo:b3:complete:8}
(a) Dini: zie \cref{lem:b3:complete:dini} --- het
overdekkingsargument. (b) Onjuist: de bewegende bult $f_n(x) =
\max(0, 1 - \abs{nx - 1})$ gaat puntsgewijs naar $0$ (voor $x > 0$
is $f_n(x) = 0$ zodra $n > 2/x$; en $f_n(0) = 0$), maar
$\norm{f_n}_\infty = 1$. (c) Onjuist: $f_n(x) = x^n$ daalt naar de
discontinue $\mathbf 1_{\{1\}}$, en $\sup_{[0,1]}\abs{f_n - f} =
\sup_{[0,1)} x^n = 1$. (d) Onjuist: $x^n \downarrow 0$
puntsgewijs op het \omterm{def:b3:topology:compact}{niet-compacte} $\intoo01$, met
$\sup_{(0,1)}x^n = 1$. Elke hypothese van Dini is dus nodig.
\end{solution}

\begin{solution}{exo:b3:complete:9}
(a) Wegens de lineariteit is $\int_0^1 fP = 0$ voor elke veelterm
$P$. Kies $P_n \to f$ uniform
(\cref{cor:b3:complete:weierstrass}): dan is $\int_0^1 f^2 =
\lim\int_0^1 fP_n = 0$, en de \omterm{def:b3:topology:continuity}{continue} $f^2 \geq 0$ met integraal
nul verdwijnt identiek.

(b) Op $\intcc01$ vormen de veeltermen in $x^2$ een algebra met de
constanten die de punten scheidt ($x \mapsto x^2$ is injectief op
$\intcc01$): \omterm{def:b3:topology:interior}{dicht} volgens Stone--Weierstrass. Op $\intcc{-1}1$
neemt $x^2$ in $\pm x$ dezelfde waarde aan, en elke veelterm in
$x^2$ ook: een uniforme limiet van zulke functies is even.
Convergeerden even functies $g_n$ uniform naar de identiteit, dan
was $x = \lim g_n(x) = \lim g_n(-x) = -x$ voor alle $x$: absurd
--- niet \omterm{def:b3:topology:interior}{dicht}. De scheidingshypothese faalt bij de paren $\{x,
-x\}$.

(c) Nee. Voor $f$ in de algebra $\mathcal A$ voortgebracht door de
constanten en $\eu^{\iu t}$ --- lineaire combinaties van
$\eu^{\iu nt}$ met $n \geq 0$ --- is $\Lambda(f) =
\int_0^{2\pi}f(t)\,\eu^{\iu t}\,\dd t = 0$ (elke
$\int_0^{2\pi}\eu^{\iu(n+1)t}\dd t = 0$). $\Lambda$ is \omterm{def:b3:topology:continuity}{continu}
voor $\norm\cdot_\infty$ ($\abs{\Lambda(f)}\leq 2\pi\norm
f_\infty$), dus verdwijnt $\Lambda$ op $\bar{\mathcal A}$; maar
$\Lambda(\eu^{-\iu t}) = 2\pi \neq 0$: dus $\eu^{-\iu t} \notin
\bar{\mathcal A}$. (Stone--Weierstrass is niet van toepassing:
$\mathcal A$ is niet gesloten onder toevoeging --- en die
belemmering is precies degene die de theorie van de holomorfe
functies in \cref{ch:b3:holomorphic} zal systematiseren.)
\end{solution}

\begin{solution}{exo:b3:complete:10}
$F_n = \bigcap_{f \in \mathcal F}\{x : \abs{f(x)} \leq n\}$ is een
doorsnede van gesloten verzamelingen: gesloten. De puntsgewijze
begrensdheid geeft $X = \bigcup_n F_n$. Baire
(\cref{thm:b3:complete:baire}) levert een $n_0$ met $U =
\mathring F_{n_0} \neq \varnothing$: op $U$ is $\abs f \leq n_0$
voor elke $f \in \mathcal F$ tegelijk.
\end{solution}

\begin{solution}{exo:b3:complete:11}
(a) Een Cauchyrij voor $\norm\cdot_{\mathcal C^1}$ is uniform
Cauchy, samen met haar afgeleiden: $f_n \to f$ en $f_n' \to g$
uniform, met $f, g$ \omterm{def:b3:topology:continuity}{continu}. Overgaan tot de limiet (de uniforme
convergentie laat dat onder de integraal toe) in $f_n(x) = f_n(0)
+ \int_0^xf_n'(t)\,\dd t$ geeft $f(x) = f(0) + \int_0^xg$: $f$ is
$\mathcal C^1$ met $f' = g$, en $\norm{f_n - f}_{\mathcal C^1} \to
0$. \omterm{def:b3:complete:complete}{Volledig}.

(b) $h(x) = \abs{x - \frac12}$ is een uniforme limiet van
$\mathcal C^1$-functies, bijvoorbeeld $h_n(x) = \sqrt{(x -
\frac12)^2 + \frac1n}$ (met $\abs{h_n - h} \leq \frac1{\sqrt n}$
via de toegevoegde grootheid), en toch is $h \notin E$: de supnorm
op $E$ is niet \omterm{def:b3:complete:complete}{volledig} --- haar \omterm{thm:b3:complete:completion}{vervollediging} is $\mathcal
C(\intcc01)$.

(c) $f_n(x) = \sin(2\pi nx)$ heeft $\norm{f_n}_\infty = 1$ en
$\norm{f_n'}_\infty = 2\pi n \to \infty$: er is geen $C$.
(Conceptueel: was het differentiëren begrensd voor de supnorm, dan
waren de twee normen uit (a)--(b) gelijkwaardig, wat $(E,
\norm\cdot_\infty)$ \omterm{def:b3:complete:complete}{volledig} zou maken --- in strijd met (b). Die
onbegrensdheid is precies waarom generieke \omterm{def:b3:topology:continuity}{continue} functies
nergens differentieerbaar kunnen zijn,
\cref{thm:b3:complete:nowherediff}.)
\end{solution}

\begin{solution}{exo:b3:complete:12}
Leg $\varepsilon > 0$ vast. Elke $F_N$ is een doorsnede over $n
\geq N$ van originelen van de gesloten $\intcc{-\varepsilon}
\varepsilon$ onder de \omterm{def:b3:topology:continuity}{continue} $x \mapsto f(nx)$: gesloten. De
hypothese zegt dat elke $x > 0$ in een zekere $F_N$ ligt. Baire,
toegepast binnen het volledige $\intcc ab$ (met willekeurige $0 <
a < b$), maakt een zekere $F_N$ \omterm{def:b3:topology:interior}{dicht} in een deelinterval; en
omdat ze gesloten is, bevat ze een interval $\intcc{a'}{b'}$ met
$0 < a' < b'$. Dan overlappen voor elke $n \geq \max(N,
\frac{a'}{b' - a'})$ de intervallen $\intcc{na'}{nb'}$ en
$\intcc{(n+1)a'}{(n+1)b'}$ elkaar (want $nb' \geq (n+1)a'$), dus
\[
\bigcup_{n \geq n_0}\intcc{na'}{nb'} \supseteq
\intco{n_0a'}{+\infty},
\]
en elke $t \geq n_0a'$ is $t = nx$ met $n \geq N$ en $x \in
\intcc{a'}{b'} \subseteq F_N$: dus $\abs{f(t)} \leq \varepsilon$.
Bijgevolg is $\limsup_{t\to\infty}\abs f \leq \varepsilon$ voor
elke $\varepsilon$: $f \to 0$. De volledige hypothese is nodig
omdat $F_N$ een heel interval aan $x$-waarden moet
\emph{overdekken} --- met alleen rationale $x$ is de vereniging
van de $F_N$ aftelbaar en geeft Baire niets; en inderdaad bestaan
er \omterm{def:b3:topology:continuity}{continue} tegenvoorbeelden die langs alle rationale stralen
verdwijnen maar niet in het oneindige.
\end{solution}

\begin{solution}{pb:b3:complete:1}
\textbf{1.} Inductie naar $k$: is $\norm{\varphi_n(t_k) - x_0}
\leq M(t_k - t_0) \leq MT \leq b$, dan ligt het punt $(t_k,
\varphi_n(t_k))$ in $R$, zodat de helling $f(t_k,
\varphi_n(t_k))$ gedefinieerd is met norm $\leq M$; en dan is voor
$t \in [t_k, t_{k+1}]$: $\norm{\varphi_n(t) - x_0} \leq
\norm{\varphi_n(t_k) - x_0} + M(t - t_k) \leq M(t - t_0) \leq MT
\leq b$.

\textbf{2.} Elk affien stuk heeft een helling van norm $\leq M$;
en een stuksgewijs affiene functie met hellingen begrensd door $M$
is $M$-lipschitz (ketting via de roosterpunten).

\textbf{3.} De familie $(\varphi_n)$ is puntsgewijs begrensd
(waarden in $\bar B(x_0, b)$) en \omterm{def:b3:complete:equicontinuous}{equicontinu} (gemeenschappelijke
lipschitzconstante $M$): Ascoli
(\cref{thm:b3:complete:ascoli}) haalt er $\varphi_{n_j} \to
\varphi$ uniform op $[t_0, t_0 + T]$ uit. De grenzen gaan over op
de limiet: $\varphi$ is $M$-lipschitz met $\varphi(t_0) = x_0$.

\textbf{4.} $R$ is \omterm{def:b3:topology:compact}{compact} en $f$ \omterm{def:b3:topology:continuity}{continu}: dus uniform \omterm{def:b3:topology:continuity}{continu}
(Heine, \cref{cor:b3:topology:heineborel}); zij
$\delta(\varepsilon)$ een modulus. Voor $t \in (t_k, t_{k+1})$
buiten het rooster is $\varphi_n'(t) = f(t_k, \varphi_n(t_k))$, en
verschillen de twee evaluatiepunten van $f$ met $\abs{t - t_k}
\leq T/n$ in de tijd en $\norm{\varphi_n(t) - \varphi_n(t_k)} \leq
MT/n$ in de ruimte. Voor $n \geq n_0$ met $(1 + M)T/n_0 <
\delta(\varepsilon)$ is dus $\norm{\Delta_n(t)} \leq
\varepsilon$.

\textbf{5.} Op elke $[t_k, t_{k+1}]$ is $\varphi_n$ affien, dus
$\varphi_n(t_{k+1}) - \varphi_n(t_k) = \int_{t_k}^{t_{k+1}}
\varphi_n'(s)\dd s$, want de afgeleide is de constante helling;
sommeren over de stukken (en het laatste bij $t$ afsnijden) geeft
$\varphi_n(t) = x_0 + \int_{t_0}^t\varphi_n'(s)\,\dd s$. Schrijven
we $\varphi_n' = f(s, \varphi_n(s)) + \Delta_n(s)$ (stuksgewijs
\omterm{def:b3:topology:continuity}{continue} integranden, met eindig veel sprongen), dan volgt de
getoonde formule.

\textbf{6.} Bij gegeven $\varepsilon$ is voor grote $j$:
$\norm{\varphi_{n_j} - \varphi}_\infty < \delta(\varepsilon)$,
dus $\norm{f(s, \varphi_{n_j}(s)) - f(s, \varphi(s))} \leq
\varepsilon$ voor alle $s$: uniforme convergentie van de
integranden, en $\int_{t_0}^t f(s,\varphi_{n_j}(s))\dd s \to
\int_{t_0}^tf(s, \varphi(s))\dd s$, uniform in $t$. Verder is
$\norm{\int_{t_0}^t\Delta_{n_j}} \leq T\sup\norm{\Delta_{n_j}} \to
0$ (vraag 4). Overgaan tot de limiet in de identiteit van vraag 5
geeft $\varphi(t) = x_0 + \int_{t_0}^tf(s, \varphi(s))\,\dd s$.

\textbf{7.} De integrand $s \mapsto f(s, \varphi(s))$ is \omterm{def:b3:topology:continuity}{continu},
dus is het rechterlid $\mathcal C^1$ in $t$ met afgeleide $f(t,
\varphi(t))$: $\varphi$ lost het beginwaardeprobleem op $[t_0, t_0
+ T]$ op. Voor de linkerhelft: stel $g(t, x) = -f(2t_0 - t, x)$,
\omterm{def:b3:topology:continuity}{continu} op de gespiegelde rechthoek met dezelfde grens $M$; een
oplossing $\psi$ van $y' = g(t,y)$ met $y(t_0) = x_0$ op $[t_0,
t_0 + T]$ levert $\varphi(t) = \psi(2t_0 - t)$, die de
oorspronkelijke vergelijking op $[t_0 - T, t_0]$ oplost; en de
twee helften verlijmen tot een $\mathcal C^1$-oplossing (beide
eenzijdige afgeleiden in $t_0$ zijn $f(t_0, x_0)$). --- Dat is de
\emph{stelling van Peano}: een \omterm{def:b3:topology:continuity}{continue} $f$ laat door elke
beginvoorwaarde een lokale oplossing toe.

\textbf{8.} De \omterm{def:b3:topology:continuity}{continuïteit} is duidelijk. Lipschitz nabij $0$
faalt: $\abs{f(h) - f(0)} = 2\sqrt h$, en $2\sqrt h \leq L h$ is
onjuist voor $h < 4/L^2$.

\textbf{9.} Voor $t \leq c$ lost $x_c \equiv 0$ op. Voor $t \geq
c$ is $x_c'(t) = 2(t - c) = 2\sqrt{(t-c)^2} = 2\sqrt{\abs{x_c}}$.
In $t = c$ zijn beide eenzijdige afgeleiden $0$: $x_c$ is
$\mathcal C^1$ en lost overal op, met $x_c(0) = 0$ voor elke $c
\geq 0$ --- samen met $x \equiv 0$ een continuüm van oplossingen
door de oorsprong.

\textbf{10.} De picarditeratie stelt $\Phi(x)(t) = x_0 + \int_0^t
f(x(s))\dd s$ op en heeft $\norm{\Phi(x) - \Phi(y)} \leq
k\norm{x - y}$ met $k < 1$ nodig op een geschikte bal --- wat
volgt uit een lipschitzgrens op $f$, onder de integraal
overgedragen. Hier laat $\sqrt\cdot$ nabij $0$ geen enkele
lipschitzgrens toe, en geen keuze van interval of bal repareert
dat. En inderdaad kon geen bewijs van eenduidigheid slagen: de
eenduidigheid is \emph{onjuist} (vraag 9).

\textbf{11.} In $c_0$ voldoen de eenheidsvectoren $e_n = (0,
\dots, 0, 1, 0, \dots)$ aan $\norm{e_n - e_m}_\infty = 1$ voor $n
\neq m$: geen enkele deelrij is Cauchy, dus is de gesloten
eenheidsbal niet \omterm{def:b3:topology:compact}{compact}. De stap die breekt is de extractie
(vraag 3): Ascoli voor $\mathcal C([t_0, t_0+T], E)$ vergt dat de
waarden leven in een ruimte waar begrensde verzamelingen relatief
\omterm{def:b3:topology:compact}{compact} zijn --- waar in $\R^d$ (Bolzano--Weierstrass), onjuist in
$c_0$; puntsgewijze extractie is er niet meer beschikbaar (en
inderdaad heeft het voorbeeld van Dieudonné helemaal geen lokale
oplossing).

\textbf{12.} \emph{Peano}: is $f$ \omterm{def:b3:topology:continuity}{continu} op een \omterm{def:b3:topology:topology}{omgeving} van
$(t_0, x_0)$ in $\R\times\R^d$, dan bestaat er een $\mathcal
C^1$-oplossing van $x' = f(t,x)$ met $x(t_0) = x_0$ op een zeker
$[t_0 - T, t_0 + T]$. \emph{Cauchy--Lipschitz}
(\cref{ch:b3:ode}): is $f$ bovendien lokaal lipschitz in de
veranderlijke $x$, dan is de oplossing \emph{uniek} (twee
oplossingen vallen samen op hun gemeenschappelijke interval) ---
bestaan én eenduidigheid. Het paar $(x' = 2\sqrt{\abs x},\ x(0) =
0)$ scheidt de twee stellingen.

\textbf{13.} $\omega(r) = Lr$: $\int_0^1\frac{\dd r}{Lr} =
+\infty$, dus voldoet het. $\omega(r) = r\log\frac1r$ (nabij $0$):
$\int\frac{\dd r}{r\log(1/r)} = \bigl[-\log\log\frac1r\bigr] \to
+\infty$ als $r \to 0$, dus voldoet het --- en toch is
$\omega(r)/r = \log\frac1r \to \infty$: niet lipschitz.
$\omega(r) = 2\sqrt r$: $\int_0^1\frac{\dd r}{2\sqrt r} =
\bigl[\sqrt r\bigr]_0^1 = 1 < \infty$, dus voldoet het niet.

\textbf{14.} Trek de twee integraalvormen van elkaar af:
\[
x_1(t) - x_2(t) = x_1(s) - x_2(s) + \int_s^t\bigl(f(v,
x_1(v)) - f(v, x_2(v))\bigr)\dd v,
\]
neem normen en gebruik de modulus van Osgood: $\delta(t) \leq
\delta(s) + \int_s^t\omega(\delta(v))\,\dd v$.

\textbf{15.} $\tau$ is goed gedefinieerd (want $\delta(t_0) = 0$)
met $\delta(\tau) = 0$ (\omterm{def:b3:topology:continuity}{continuïteit}) en $\delta > 0$ op
$\intoc\tau{t_1}$. Leg $s \in \intoo\tau{t_1}$ vast. Dan is $u$
$\mathcal C^1$ met $u(s) = \delta(s) > 0$, is $u \geq \delta$ op
$[s, t_1]$ (vraag 14), en is $u' = \omega(\delta) \leq \omega(u)$
($\omega$ is niet-dalend en $u > 0$). Delen en integreren geeft
\[
\int_{\delta(s)}^{u(t_1)}\frac{\dd r}{\omega(r)}
= \int_s^{t_1}\frac{u'(v)}{\omega(u(v))}\,\dd v \leq t_1 - s
\leq t_1 - \tau .
\]
Hoewel $u$ van $s$ afhangt, geldt $u(t_1) \geq \delta(t_1)$ voor
elke $s$, zodat het linkerlid minstens
$\int_{\delta(s)}^{\delta(t_1)}\frac{\dd r}{\omega(r)}$ is, en dat
gaat naar $+\infty$ als $s \downarrow \tau$ (dan is $\delta(s) \to
\delta(\tau) = 0$, en de integraal divergeert in $0$): het
begrensde rechterlid wordt tegengesproken. Bijgevolg is $\delta
\equiv 0$: eenduidigheid.

\textbf{16.} Lipschitz is $\omega = Lr$: de eenduidigheid is
teruggevonden. Voor $x' = x\log\frac1{\abs x}$ voldoet het
rechterlid nabij $0$ aan de Osgood-modulus $\omega(r) =
r\log\frac1r$ (middelwaardeongelijkheid op $x \mapsto
x\log\frac1x$, waarvan de afgeleide $\log\frac1x - 1$ onbegrensd
is --- Lipschitz faalt, Osgood geldt): dus eenduidige oplossingen;
en merk op dat $x \equiv 0$ er een is, zodat geen andere oplossing
$0$ kan raken. Voor $\omega(r) = 2\sqrt r$ is
$\int_0^{h^2}\frac{\dd r}{2\sqrt r} = h$ --- het
``Osgood-budget'' om van $0$ naar hoogte $h^2$ te klimmen is
precies de tijd $h$, en inderdaad is $x_c(c + h) = h^2$: de
oplossing besteedt precies de tijd $h$ aan wat de convergente
integraal toestaat. Divergentie van de integraal is de
onmogelijkheid om $0$ in eindige tijd te verlaten; convergentie is
de ontsnappingsroute.

\textbf{17.} $G(t) = \int_{t_0}^t(Le(s) + \eta(s))\dd s$ is
$\mathcal C^1$ met $G' = Le + \eta \leq LG + \sup\eta$ (wegens de
hypothese $e \leq G$). Dan is $\bigl(\eu^{-Lt}(G +
\tfrac{\sup\eta}L)\bigr)' = \eu^{-Lt}(G' - LG - \sup\eta) \leq 0$:
de uitdrukking tussen haakjes daalt, dus $G(t) + \frac{\sup\eta}L
\leq \eu^{L(t - t_0)}\bigl(G(t_0) + \frac{\sup\eta}L\bigr) =
\eu^{L(t-t_0)}\frac{\sup\eta}L$, oftewel $e(t) \leq G(t) \leq
\frac{\sup\eta}L(\eu^{L(t-t_0)} - 1)$.

\textbf{18.} Kwantitatief defect: op $(t_k, t_{k+1})$ is
$\Delta_n(t) = f(t_k, \varphi_n(t_k)) - f(t, \varphi_n(t))$ met
$\abs{t - t_k} \leq T/n$ en $\norm{\varphi_n(t) -
\varphi_n(t_k)} \leq MT/n$, dus $\norm{\Delta_n} \leq L'T/n +
LMT/n$. Trek de integraalidentiteiten voor $\varphi_n$ (vraag 5)
en $\varphi$ van elkaar af en gebruik de lipschitzgrens:
\[
\norm{\varphi_n(t) - \varphi(t)} \leq
\int_{t_0}^t L\,\norm{\varphi_n - \varphi}(s)\,\dd s +
\int_{t_0}^t\norm{\Delta_n(s)}\,\dd s,
\]
en vraag 17 met $\eta = \norm{\Delta_n}$ geeft de aangegeven
grens $O(1/n)$, waarbij $\varphi$ uniek is volgens
Cauchy--Lipschitz (of Osgood). Ook zonder snelheden: elke deelrij
van de equibegrensde en equi-lipschitz rij $(\varphi_n)$ heeft een
sub-deelrij die (Ascoli plus Deel II) naar \emph{een} oplossing
convergeert, die de eenduidigheid dwingt gelijk te zijn aan
$\varphi$; en een rij waarvan elke deelrij een sub-deelrij met
dezelfde limiet heeft, convergeert.

\textbf{19.} Uit $\varphi_n(t_k) = 0$ volgt met de helling
$2\sqrt0 = 0$ dat $\varphi_n(t_{k+1}) = 0$; per inductie is
$\varphi_n \equiv 0$, convergerend naar de nuloplossing. Vanaf
$\varepsilon > 0$: op $[\varepsilon', \infty)$ met $\varepsilon' <
\varepsilon$ is de functie $2\sqrt x$ lipschitz, dus is vraag 18
van toepassing en convergeert Euler naar de unieke oplossing door
$(0, \varepsilon)$, namelijk $x(t) = (t + \sqrt\varepsilon)^2$
(controle: $x' = 2(t + \sqrt \varepsilon) = 2\sqrt x$). Als
$\varepsilon \to 0$ gaat $(t + \sqrt\varepsilon)^2 \to t^2$
uniform op $[0,1]$: de dubbele limiet landt op $x_0(t) = t^2$ en
niet op $0$. Een willekeurig kleine verstoring van het beginpunt
stuurt het schema van de ene oplossing naar de andere:
niet-eenduidigheid gelezen als numerieke instabiliteit.

\textbf{20.} Niet leeg: Deel II. Elke oplossing voldoet aan
$\norm{x'} = \norm{f(t, x)} \leq M$: $\mathcal S$ is uniform
$M$-lipschitz en uniform begrensd (waarden in $\bar B(x_0, b)$).
Gesloten: convergeert $x_n \in \mathcal S$ uniform naar $x$, ga
dan tot de limiet over in $x_n(t) = x_0 + \int_{t_0}^tf(s,
x_n(s))\dd s$ (de integranden convergeren uniform wegens de
uniforme \omterm{def:b3:topology:continuity}{continuïteit} van $f$ op de \omterm{def:b3:topology:compact}{compacte} $R$): dus $x \in
\mathcal S$. Ascoli: $\mathcal S$ is een gesloten, begrensde,
\omterm{def:b3:complete:equicontinuous}{equicontinue} deelverzameling van $\mathcal C([t_0, t_0+T],
\R^d)$: \omterm{def:b3:topology:compact}{compact}.

\textbf{21.} Elke oplossing is niet-dalend ($x' = 2\sqrt{\abs x}
\geq 0$) met $x(0) = 0$, dus $x \geq 0$. Stel $c = \sup\{t \in
[0,1] : x(t) = 0\}$ (eventueel $c = \infty$ als $x \equiv 0$, en
dan is $x = x_\infty$). Voor $t > c$ is $x > 0$ (monotonie plus de
definitie van $c$), en daar is $(\sqrt x)' = \frac{x'}{2\sqrt x} =
1$, dus $\sqrt{x(t)} = t - c$ (\omterm{def:b3:topology:continuity}{continuïteit} in $c$): $x = x_c$.
\omterm{def:b3:topology:continuity}{Continuïteit} van $c \mapsto x_c$: voor $c, c' \in [0, 1]$ is
$\sup_{[0,1]}\abs{(t - c)_+^2 - (t - c')_+^2} \leq 2\abs{c - c'}$
(de afbeelding $c \mapsto (t-c)_+^2$ is $2$-lipschitz, uniform in
$t \in [0,1]$), en omdat $x_c \equiv 0$ op $[0,1]$ voor elke $c
\geq 1$, herleidt de familie zich tot $\mathcal S = \{x_c : c \in
[0,1]\}$ (met $x_1 = 0 = x_\infty$). Dus is $\mathcal S$ het beeld
van het \omterm{def:b3:topology:compact}{compacte} \omterm{def:b3:topology:connected}{samenhangende} $[0,1]$ onder de \omterm{def:b3:topology:continuity}{continue
afbeelding} $c \mapsto x_c$: \omterm{def:b3:topology:compact}{compact} en \omterm{def:b3:topology:connected}{samenhangend}. De trechter
van oplossingen is een \omterm{def:b3:topology:continuity}{continu} lijnstuk dat van $t^2$ (onmiddellijke
ontsnapping) tot $0$ (eeuwige rust) loopt.

\textbf{22.} De evaluatie $\operatorname{ev}_t\colon \mathcal
C([0,1]) \to \R$, $x \mapsto x(t)$, is \omterm{def:b3:topology:continuity}{continu} ($\abs{x(t) -
y(t)} \leq \norm{x - y}_\infty$), dus is $\mathcal S(t) =
\operatorname{ev}_t(\mathcal S)$ een \omterm{def:b3:topology:continuity}{continu} beeld van een
\omterm{def:b3:topology:compact}{compacte} verzameling: \omterm{def:b3:topology:compact}{compact} --- en van een \omterm{def:b3:topology:connected}{samenhangende}:
\omterm{def:b3:topology:connected}{samenhangend}. Voor het voorbeeld doorloopt $x_c(t) = (t - c)_+^2$,
terwijl $c$ het interval $[0, 1]$ doorloopt, \omterm{def:b3:topology:continuity}{continu} alle waarden
van $t^2$ (bij $c = 0$) tot $0$ (bij $c \geq t$): dus $\mathcal
S(t) = \intcc0{t^2}$. Op elk tijdstip is de doorsnede van de
trechter een vol lijnstuk: tussen rust en maximale ontsnapping
wordt elk compromis door een echte oplossing gerealiseerd.

\textbf{23.} Trek de integraalvormen $x(t) = x_0 +
\int_{t_0}^tf(s, x(s))\dd s$ en haar tegenhanger voor $y$ van
elkaar af, en stel $e(t) = \norm{x(t) - y(t)}$ en $d = \norm{x_0 -
y_0}$:
\[
e(t) \leq d + \int_{t_0}^tL\,e(s)\dd s = G(t).
\]
Dan is $G(t_0) = d$ en $G' = Le \leq LG$, dus $\bigl(\eu^{-L(t -
t_0)}G\bigr)' \leq 0$ en $G(t) \leq d\,\eu^{L(t - t_0)}$; en
bijgevolg $e(t) \leq d\,\eu^{L(t-t_0)}$. Scherpte: voor $x' = Lx$
zijn de oplossingen door $x_0$ en $y_0$ gelijk aan
$x_0\eu^{L(t - t_0)}$ en $y_0\eu^{L(t-t_0)}$, met afstand precies
$d\,\eu^{L(t-t_0)}$. Met $d = 0$ is $e \equiv 0$: de eenduidigheid
van Cauchy--Lipschitz, in twee regels opnieuw afgeleid. En voor
vaste $t \in [t_0, t_0 + T]$ is $\norm{x(t; x_0) - x(t; y_0)} \leq
\eu^{LT}\norm{x_0 - y_0}$: de stroming is lipschitz in de
beginvoorwaarde --- deterministische afhankelijkheid, tegen een
beheerste exponentiële prijs.

\textbf{24.} Op $\intoc01$ is $\Omega$ goed gedefinieerd (de
integraal convergeert in $0$ per hypothese) en $\mathcal C^1$ met
$\Omega' = \frac1\omega > 0$: een stijgende bijectie op
$\intoc0{\Omega(1)}$, met $\Omega(x) \to 0$ als $x \to 0^+$. Haar
inverse $g \colon \intoc0{\Omega(1)} \to \intoc01$ is $\mathcal
C^1$ met
\[
g'(t) = \frac1{\Omega'(g(t))} = \omega\bigl(g(t)\bigr) > 0,
\qquad g(t) \xrightarrow[t \to 0^+]{} 0 .
\]
Zet $g$ voort met $0$ voor $t \leq 0$: de \omterm{def:b3:topology:continuity}{continuïteit} is
duidelijk, en in $t = 0$ is voor $t > 0$
\[
\frac{g(t)}t = \frac1t\int_0^tg'(s)\dd s
= \frac1t\int_0^t\omega\bigl(g(s)\bigr)\dd s
\leq \omega\bigl(g(t)\bigr) \xrightarrow[t\to0^+]{} 0
\]
($\omega$ is niet-dalend, $g$ stijgend, en $\omega(0^+) = 0$):
dus $g'(0) = 0 = f(g(0))$, en $g' = f(g)$ geldt aan beide kanten
van $0$. Bijgevolg zijn $x \equiv 0$ en $g$ twee verschillende
$\mathcal C^1$-oplossingen door $(0, 0)$: convergeert
$\int_0\frac{\dd r}{\omega(r)}$, dan faalt de eenduidigheid ---
de divergentie uit vraag 15 is precies de grens. Voor $\omega(r) =
2\sqrt r$ is $\Omega(x) = \int_0^x\frac{\dd r}{2\sqrt r} = \sqrt
x$ en $g(t) = t^2$, en tijdverschuivingen geven de hele familie
$x_c$ uit Deel III.

\textbf{25.} Met stap $h = \frac1n$ is $\varphi_n(t_{k+1}) =
\varphi_n(t_k) + h\,\varphi_n(t_k) = (1 + h)\varphi_n(t_k)$, dus
$\varphi_n(1) = (1 + \frac1n)^n$ na $n$ stappen. Ontwikkeling:
\[
n\log\Bigl(1 + \frac1n\Bigr)
= n\Bigl(\frac1n - \frac1{2n^2} + O\Bigl(\frac1{n^3}\Bigr)
\Bigr) = 1 - \frac1{2n} + O\Bigl(\frac1{n^2}\Bigr),
\]
en na exponentiëren: $(1 + \frac1n)^n = \eu\,\eu^{-1/(2n) +
O(n^{-2})} = \eu\bigl(1 - \frac1{2n} + O(n^{-2})\bigr)$. De fout
in $t = 1$ is dus $\eu - \varphi_n(1) = \frac{\eu}{2n} +
O(n^{-2})$: de $O(\frac1n)$ van vraag 18, hier met haar exacte
constante $\frac\eu2$. Numeriek voor $n = 10$: $1.1^{10} =
2.5937424601$ ($1.1^2 = 1.21$, $1.1^4 = 1.4641$, $1.1^8 =
2.14358881$, maal $1.21$), en $\eu - 2.59374 \approx 0.12454$
tegenover de asymptotische voorspelling $\frac{\eu}{20} \approx
0.13591$: overeenstemming tot op de correctie $O(n^{-2})$, waarvan
de hoofdterm hier de voorspelling naar de waargenomen waarde toe
verlaagt.
\end{solution}
